\section{Related Work}
\label{sec:related}

Predictive-maintenance models often begin with supervised classification, reconstruction error, distance-based anomaly detection, or streaming thresholds. These methods are useful engineering baselines, but their common output is a timestamped score rather than an episode-level maintenance decision. EVT-based streaming detectors such as SPOT use peaks-over-threshold modelling to adapt anomaly thresholds in streams \citep{Siffer2017SPOT}; they address threshold calibration, but they do not by themselves resolve operating-regime mixture, duplicate alarms per failure, or the distinction between clustered observations and independent maintenance events.

Classical declustering and extremal-index estimation address dependence among extremes. Ferro and Segers proposed an inference scheme for clusters of extreme values using inter-exceedance times and the extremal index \citep{FerroSegers2003Clusters}. In dynamical systems, hitting-time statistics connect returns to target neighborhoods with extreme-value laws \citep{Freitas2008Hitting}. Periodic points can generate an extremal index smaller than one, giving a mathematical route from recurrence to clustered extremes \citep{Freitas2012Periodicity}. The finite-sample industrial setting is less clean: noise may make deterministic recurrence appear closer to independence, and observed clustering can also arise from smoothing, missingness, persistence, or regime mixture \citep{Faranda2013Noise}. Recent multivariate work extends dynamical EVT to cross-sectional dependence among vector-valued observables \citep{Aimino2024Multivariate}, which motivates the lag-pattern diagnostics used here.

The industrial datasets used by the repository are public predictive-maintenance resources with different independence units. MetroPT is an air-production-unit telemetry dataset for predictive maintenance \citep{Veloso2022MetroPT}. MetroPT2 is a replication-style benchmark released through Zenodo \citep{MetroPT2Zenodo2023}. SCANIA Component X is a fleet-level multivariate time-series dataset with operational histories and repair records for an anonymized component \citep{Kharazian2025Scania}. These datasets support useful case studies, but they do not remove the need for event-level accounting: one long fault interval is not equivalent to thousands of independent statistical units. The methodological gap addressed here is the absence of a reproducible layer that connects dynamical extreme diagnostics to alarm episodes, sensitivity analyses, and claim constraints.
